problem:  
Let \( G \) be a compact connected semisimple Lie group, \( P \subset G \) a parabolic subgroup, and \( M = G/P \) its associated generalized flag manifold. Denote by  
\[
\mathfrak{g} = \mathfrak{p} \oplus \mathfrak{m} = \mathfrak{p} \oplus \bigoplus_{i=1}^r \mathfrak{m}_i
\]
the isotropy decomposition into irreducible \( \operatorname{Ad}(P) \)-modules, where \( r \) is the number of isotropy summands.  
Let \( g_{KE} \) be the canonical invariant Kähler–Einstein metric on \( M \).  

Each \( G \)-invariant Riemannian metric on \( M \) corresponds to a point \( (x_1, \dots, x_r) \in (\mathbb{R}_{>0})^r \), scaling the components in the decomposition. The Einstein condition then gives a system of algebraic equations
\[
F_i(x_1, \dots, x_r) = \lambda x_i, \qquad i = 1, \dots, r.
\]

**Problem:**  
At the point corresponding to \( g_{KE} \), consider the linearization of the system \( F_i(x) = \lambda x_i \), i.e. the Jacobian matrix \(J(g_{KE})\).  

Determine whether there exist flag manifolds \(G/P\) for which \(J(g_{KE})\) is *singular*, meaning that the Kähler–Einstein solution is a degenerate critical point of the homogeneous Einstein equations. In such a case, would there exist a continuous family of invariant Einstein metrics (Kähler or non‑Kähler) branching from \(g_{KE}\)?  

Equivalently: **For which flag manifolds is the canonical Kähler–Einstein metric infinitesimally deformable (within invariant metrics) as an Einstein metric?**

Why is it a "good" problem:  
This formulation turns a loose topological question into a precise structural one about *linearized Einstein stability in the invariant setting*, connecting algebraic representation theory with geometric analysis. The problem digs into the interface between the algebraic parametrization of invariant metrics and the functional-analytic properties of the Einstein operator.  

Although homogeneous Einstein metrics on many flag manifolds are known to be isolated, it is not known in full generality when—and why—the linearized Einstein operator degenerates at the Kähler–Einstein point. This degeneracy would correspond to new, possibly non‑Kähler, invariant Einstein directions, revealing interactions between the representation‑theoretic structure of the isotropy representation and the analytic rigidity of the Einstein condition.  

In essence, the question asks whether algebraic degeneracy reflects geometric deformability—a concrete, conceptually rich bridge between algebraic and differential perspectives on Einstein geometry. It remains nontrivial, specific, and open, while being logically well‑posed and grounded in the actual structure of homogeneous spaces.